% !TeX root = closed-systems-from-comparison-completeness.tex
% ============================================================
\chapter*{Orientation to Part III}
\phantomsection
\label{part:intro-transport}
% ============================================================
\begin{chapterorientation}
\orientationitem{Settled}{The physical quotient and the subsystem-attribution obstruction are fixed by Part~II.}
\orientationitem{Task}{Part~III classifies all admissible representative enrichment and identifies the first transport obstruction.}
\orientationitem{Handoff}{The output is the triangle--loop rigidity bridge used by the observer, spacetime, and curvature chapters.}
\end{chapterorientation}

Part~III studies what remains after quotient semantics has been fixed.
Representative information can still enter a calculation, but only
through controlled loci: representative choice and morphism-level
transport.  The part proves that these loci exhaust admissible
enrichment.

\Cref{ch:locus,ch:lift} give the action-groupoid and lift calculus.
\Cref{ch:triangles} isolates the first finite multi-point witness of
transport obstruction, while \cref{ch:descent_obstruction} gives the
global loop representation.  \Cref{ch:bridge} proves that the triangle
and loop forms are the same obstruction datum.  This identification is
the algebraic hinge on which the later geometry turns.
